\chapter{Order Statistics}\label{ch:iv1:order-statistics}

\begin{ivtrains}
Working with the sorted values of a sample: maximum and median, spacings, the highest serial number seen, the second-highest bid, records. Ranks are the common thread, because the rank pattern of a sample does not depend on its distribution.
\end{ivtrains}

\begin{question}[$\star$ \iqroles{QR, QT, Trader, Analyst} \iqfirm{any} \iqlevel{junior}]\label{iq:iv1:order-statistics:1}
You draw $n$ independent uniform numbers on $[0, 1]$. What are the averages of the largest, of the smallest and of the range? Give the probability that the largest is below $0.9$ when $n = 10$.
\end{question}

\begin{question}[$\star$ \iqroles{QR, QT, DS, Analyst} \iqfirm{any} \iqlevel{junior}]\label{iq:iv1:order-statistics:2}
You draw three independent uniform numbers on $[0, 1]$. What is the density of their median? What are its mean and variance? Which estimates the centre of the distribution better, the median of the three or their average?
\end{question}

\begin{question}[$\star$ \iqroles{QR, QT, Trader, DS} \iqfirm{any} \iqlevel{junior}]\label{iq:iv1:order-statistics:3}
Desk A sends 10 orders and desk B sends 20, and all 30 latencies are independent draws from the same continuous distribution. What is the probability that the fastest order belongs to desk B? That the three fastest all belong to desk B? Does the answer depend on the distribution?
\end{question}

\begin{question}[$\star\star$ \iqroles{QR, QT, DS, Analyst} \iqfirm{any} \iqlevel{junior}]\label{iq:iv1:order-statistics:4}
A batch of machines carries the serial numbers $1, 2, \dots, N$, with $N$ unknown. You see five different machines, picked at random, and the highest serial number among them is 60. Give an unbiased estimate of $N$, and derive it. Why is twice the sample mean minus one a worse estimator, even though it is unbiased as well?
\end{question}

\begin{question}[$\star\star$ \iqroles{QR, QT, DS} \iqfirm{any} \iqlevel{junior}]\label{iq:iv1:order-statistics:5}
You have $n$ independent draws from a uniform on $[0, \theta]$. Compare the variance of $2\bar{X}$ with the variance of $\frac{n + 1}{n}\max X_i$ for $n = 10$. Both are unbiased. Which multiple of the maximum has the smallest mean squared error, and what is it worth?
\end{question}

\begin{question}[$\star\star$ \iqroles{QR, QT, Trader, Analyst} \iqfirm{market maker, hedge fund} \iqlevel{junior}]\label{iq:iv1:order-statistics:6}
Five bids are independent and uniform on $[0, 1]$. What is the average of the second-highest bid? If you are told that the highest bid is $0.8$, what is the average of the second-highest then? Separately, you roll four fair dice: what is the average of the second-highest number shown?
\end{question}

\begin{question}[$\star\star$ \iqroles{QR, QT, DS, Risk} \iqfirm{any} \iqlevel{junior}]\label{iq:iv1:order-statistics:7}
For a large sample from a normal distribution, how much more variable is the sample median than the sample mean? What is the same comparison for the Laplace distribution, with density $\frac12 e^{-|x|}$? What does this suggest for fat-tailed returns?
\end{question}

\begin{question}[$\star\star$ \iqroles{QR, QT, Risk, QD} \iqfirm{any} \iqlevel{junior}]\label{iq:iv1:order-statistics:8}
Five servers fail after independent exponential times with mean one year. When, on average, does the first fail? The third? The last? Are the time of the first failure and the time from the first failure to the last independent?
\end{question}

\begin{question}[$\star\star$ \iqroles{QR, QT, DS} \iqfirm{any} \iqlevel{junior}]\label{iq:iv1:order-statistics:9}
You draw independent numbers from a continuous distribution until one beats your first draw. Let $N$ be the number of further draws you need. What is $P(N > n)$? What is the smallest $n$ with $P(N \le n) \ge \frac12$, and what is the mean of $N$? Are the events ``draw 5 is a record'' and ``draw 10 is a record'' independent?
\end{question}

\begin{question}[$\star\star$ \iqroles{QR, QT, DS, Risk} \iqfirm{any} \iqlevel{junior}]\label{iq:iv1:order-statistics:10}
You have 10 independent draws from an unknown continuous distribution. What is the probability that its true median lies between the smallest and the largest draw? Between the third smallest and the third largest? Why can you say this with no knowledge of the distribution?
\end{question}

\begin{question}[$\star\star\star$ \iqroles{QR, QT, DS} \iqfirm{hedge fund, market maker} \iqlevel{junior, experienced}]\label{iq:iv1:order-statistics:11}
A researcher backtests 1000 independent strategies, all with a true Sharpe ratio of 0. Each measured annual Sharpe ratio is normal, centred on 0, with a standard deviation of $1/3$, which corresponds to 9 years of data. About what is the best measured Sharpe ratio, working without a table? How does the expected best grow with the number of strategies tried?
\end{question}

\begin{question}[$\star\star\star$ \iqroles{QR, QT, Risk, DS} \iqfirm{any} \iqlevel{junior, experienced}]\label{iq:iv1:order-statistics:12}
You draw $n$ independent uniform numbers on $[0, 1]$. Find the distribution of the range. How large must $n$ be for the probability to be at least 95\% that the range covers at least 90\% of the interval? What can you then say, with the same $n$, about any continuous distribution?
\end{question}

\begin{question}[$\star\star\star$ \iqroles{QR, QT, DS} \iqfirm{hedge fund} \iqlevel{junior, experienced}]\label{iq:iv1:order-statistics:13}
A sample of size 11 comes from a uniform on $[\theta, \theta + 1]$, with $\theta$ unknown. Compare three estimators of the centre: the sample mean, the sample median and the midrange, which is the average of the smallest and largest. Give the three variances. Why does the midrange win here, and at what rate does its variance fall as the sample grows?
\end{question}

\begin{question}[$\star\star\star$ \iqroles{QR, QT, LLE, QD} \iqfirm{market maker, crypto} \iqlevel{junior, experienced}]\label{iq:iv1:order-statistics:14}
Ten orders arrive at independent uniform times within one second. What is the probability that no two of them are within 10 milliseconds of each other? What is the average of the smallest gap between neighbouring orders?
\end{question}

\begin{question}[$\star\star\star\star$ \iqroles{QR, QT} \iqfirm{hedge fund, market maker} \iqlevel{experienced}]\label{iq:iv1:order-statistics:15}
Find, in closed form, the average of the largest of three independent standard normal variables, and then the average of the largest of four. Your method should not need numerical integration. How far does the same method go, and what stops it?
\end{question}
